\documentclass[10pt,a4paper]{amsart}
\usepackage{amsmath,amssymb,amsthm,graphicx}
\usepackage{verbatim}
\newenvironment{reference}{\comment}{\endcomment}
%\newenvironment{reference}{}{}
\newenvironment{slogan}{\comment}{\endcomment}
\newenvironment{history}{\comment}{\endcomment}
% For commutative diagrams we use Xy-pic
\usepackage[all]{xy}

% We use 2cell for 2-commutative diagrams.
\xyoption{2cell}
% Theorem environments.
%
\theoremstyle{plain}
\newtheorem{theorem}[subsection]{Theorem}
\newtheorem{proposition}[subsection]{Proposition}
\newtheorem{lemma}[subsection]{Lemma}

\theoremstyle{definition}
\newtheorem{definition}[subsection]{Definition}
\newtheorem{example}[subsection]{Example}
\newtheorem{exercise}[subsection]{Exercise}
\newtheorem{situation}[subsection]{Situation}

\theoremstyle{remark}
\newtheorem{remark}[subsection]{Remark}
\newtheorem{remarks}[subsection]{Remarks}

\numberwithin{equation}{subsection}

% Macros
%
\def\lim{\mathop{\mathrm{lim}}\nolimits}
\def\colim{\mathop{\mathrm{colim}}\nolimits}
\def\Spec{\mathop{\mathrm{Spec}}}
\def\Hom{\mathop{\mathrm{Hom}}\nolimits}
\def\Ext{\mathop{\mathrm{Ext}}\nolimits}
\def\SheafHom{\mathop{\mathcal{H}\!\mathit{om}}\nolimits}
\def\SheafExt{\mathop{\mathcal{E}\!\mathit{xt}}\nolimits}
\def\Sch{\mathit{Sch}}
\def\Mor{\mathop{\mathrm{Mor}}\nolimits}
\def\Ob{\mathop{\mathrm{Ob}}\nolimits}
\def\Sh{\mathop{\mathit{Sh}}\nolimits}
\def\NL{\mathop{N\!L}\nolimits}
\def\CH{\mathop{\mathrm{CH}}\nolimits}
\def\proetale{{pro\text{-}\acute{e}tale}}
\def\etale{{\acute{e}tale}}
\def\QCoh{\mathit{QCoh}}
\def\Ker{\mathop{\mathrm{Ker}}}
\def\Im{\mathop{\mathrm{Im}}}
\def\Coker{\mathop{\mathrm{Coker}}}
\def\Coim{\mathop{\mathrm{Coim}}}

% Boxtimes
%
\DeclareMathSymbol{\boxtimes}{\mathbin}{AMSa}{"02}

%
% Macros for moduli stacks/spaces
%
\def\QCohstack{\mathcal{QC}\!\mathit{oh}}
\def\Cohstack{\mathcal{C}\!\mathit{oh}}
\def\Spacesstack{\mathcal{S}\!\mathit{paces}}
\def\Quotfunctor{\mathrm{Quot}}
\def\Hilbfunctor{\mathrm{Hilb}}
\def\Curvesstack{\mathcal{C}\!\mathit{urves}}
\def\Polarizedstack{\mathcal{P}\!\mathit{olarized}}
\def\Complexesstack{\mathcal{C}\!\mathit{omplexes}}
% \Pic is the operator that assigns to X its picard group, usage \Pic(X)
% \Picardstack_{X/B} denotes the Picard stack of X over B
% \Picardfunctor_{X/B} denotes the Picard functor of X over B
\def\Pic{\mathop{\mathrm{Pic}}\nolimits}
\def\Picardstack{\mathcal{P}\!\mathit{ic}}
\def\Picardfunctor{\mathrm{Pic}}
\def\Deformationcategory{\mathcal{D}\!\mathit{ef}}
\title{Vanishing of the local cohomology of the punctured closed set makes the comparison map between two local cohomology modules surjective, injective or bijective}
\author{The Stacks Project Authors}
\date{stacks-project:local-cohomology.tex:local-snapshot-20261012; GNU FDL 1.2 or later}
\begin{document}
\maketitle
\noindent\textit{Editorial note.} All delivered mathematical text is the author's own, copied from the pinned Stacks Project chapter source sources/stacks-local-cohomology/pinned-original-source.tex; no statement or proof step was rewritten, reordered or removed.\par\medskip
\noindent\textit{Editorial note.} Preamble provenance: the Stacks Project keeps one single shared preamble for every chapter; this card takes it directly from the staged sources/stacks-local-cohomology/preamble.tex as three verbatim spans (lines 8-12, 14-18 and 154-217), so no derived file was written and nothing was added or redefined. The wrapper class and the base packages are supplied by the assembler.\par\medskip
\noindent\textit{Editorial note.} Bibliography provenance: neither the counted statement nor its proof contains a citation (this lane rescans the whole span before writing the card), so the card carries no bibliography and no entry was reconstructed.\par\medskip
\noindent\textit{Editorial note.} Selection provenance: this card was opened under the file-based control contract (CONTROL.md), with the mechanical pre-card gate run on its statement and proof interval before the card was written.\par\medskip
\noindent\textit{Allowed standard prerequisites (editorial).}\begin{itemize}
\item Local cohomology H\textasciicircum{}i\_\{V(I)\}(M) of a module with support in the closed subset V(I), and its localisation formula H\textasciicircum{}i\_\{V(I)\}(M)\_\textbackslash{}mathfrak p = H\textasciicircum{}i\_\{\textbackslash{}mathfrak pA\_\textbackslash{}mathfrak p\}(M\_\textbackslash{}mathfrak p). The author's section on generalities about local cohomology, in which the counted statement occurs. Supplies the two modules compared by the counted map and the local vanishing hypothesis.
\item The lemma that elements of a local cohomology module come from a strictly bigger finitely generated ideal. Author lemma lemma-elements-come-from-bigger, whose statement is reproduced verbatim as a context span of this card (source lines 451-470). Provides the step that increases the auxiliary ideal in both halves of the proof and thereby contradicts maximality.
\end{itemize}\medskip
\begin{lemma}
\label{lemma-isomorphism}
Let $I \subset I' \subset A$ be finitely generated ideals of a
Noetherian ring $A$. Let $M$ be an $A$-module. Let $i \geq 0$ be an integer.
Consider the map
$$
\Psi : H^i_{V(I')}(M) \to H^i_{V(I)}(M)
$$
The following are true:
\begin{enumerate}
\item if $H^i_{\mathfrak pA_\mathfrak p}(M_\mathfrak p) = 0$
for all $\mathfrak p \in V(I) \setminus V(I')$, then
$\Psi$ is surjective,
\item if $H^{i - 1}_{\mathfrak pA_\mathfrak p}(M_\mathfrak p) = 0$
for all $\mathfrak p \in V(I) \setminus V(I')$, then
$\Psi$ is injective,
\item if $H^i_{\mathfrak pA_\mathfrak p}(M_\mathfrak p) =
H^{i - 1}_{\mathfrak pA_\mathfrak p}(M_\mathfrak p) = 0$
for all $\mathfrak p \in V(I) \setminus V(I')$, then
$\Psi$ is an isomorphism.
\end{enumerate}
\end{lemma}

\begin{lemma}
\label{lemma-elements-come-from-bigger}
Let $I \subset A$ be a finitely generated ideal of a ring $A$.
Let $\mathfrak p$ be a prime ideal. Let $M$ be an $A$-module.
Let $i \geq 0$ be an integer and consider the map
$$
\Psi :
\colim_{f \in A, f \not \in \mathfrak p} H^i_{V((I, f))}(M)
\longrightarrow
H^i_{V(I)}(M)
$$
Then
\begin{enumerate}
\item $\Im(\Psi)$ is the set of elements which map to zero in
$H^i_{V(I)}(M)_\mathfrak p$,
\item if $H^{i - 1}_{V(I)}(M)_\mathfrak p = 0$, then $\Psi$ is injective,
\item if $H^{i - 1}_{V(I)}(M)_\mathfrak p = H^i_{V(I)}(M)_\mathfrak p = 0$,
then $\Psi$ is an isomorphism.
\end{enumerate}
\end{lemma}

\begin{proof}
Proof of (1).
Let $\xi \in H^i_{V(I)}(M)$. Since $A$ is Noetherian, there exists a
largest ideal $I \subset I'' \subset I'$ such that $\xi$ is the image
of some $\xi'' \in H^i_{V(I'')}(M)$. If $V(I'') = V(I')$, then we are
done. If not, choose a generic point $\mathfrak p \in V(I'')$ not in $V(I')$.
Then we have $H^i_{V(I'')}(M)_\mathfrak p =
H^i_{\mathfrak pA_\mathfrak p}(M_\mathfrak p) = 0$ by assumption.
By Lemma \ref{lemma-elements-come-from-bigger} we can increase $I''$
which contradicts maximality.

\medskip\noindent
Proof of (2). Let $\xi' \in H^i_{V(I')}(M)$ be in the kernel of $\Psi$.
Since $A$ is Noetherian, there exists a
largest ideal $I \subset I'' \subset I'$ such that $\xi'$
maps to zero in $H^i_{V(I'')}(M)$. If $V(I'') = V(I')$, then we are
done. If not, then choose a generic point $\mathfrak p  \in V(I'')$
not in $V(I')$. Then we have $H^{i - 1}_{V(I'')}(M)_\mathfrak p =
H^{i - 1}_{\mathfrak pA_\mathfrak p}(M_\mathfrak p) = 0$ by assumption.
By Lemma \ref{lemma-elements-come-from-bigger} we can increase $I''$
which contradicts maximality.

\medskip\noindent
Part (3) is formal from parts (1) and (2).
\end{proof}

\end{document}
